\documentclass[11pt]{article}
\usepackage[a4paper,margin=1.9cm]{geometry}
\usepackage[T1]{fontenc}
\usepackage{textcomp}
\usepackage[spanish]{babel}
\usepackage{amsmath,amssymb}
\usepackage{enumitem}
\usepackage{tikz}
\usetikzlibrary{calc,arrows.meta,patterns,decorations.pathmorphing}
\usepackage{xcolor}
\usepackage{multicol}
\usepackage{parskip}
\usepackage{lmodern}
\renewcommand{\familydefault}{\sfdefault}

\definecolor{unalblue}{RGB}{31,73,124}

\newlist{opciones}{enumerate}{1}
\setlist[opciones]{label=(\alph*),itemsep=1pt,topsep=2pt,leftmargin=1.8em,font=\normalfont}

\newcommand{\vect}[2]{\begin{pmatrix}#1\\#2\end{pmatrix}}
\newcommand{\vectiii}[3]{\begin{pmatrix}#1\\#2\\#3\end{pmatrix}}

\newcommand{\examheader}{%
  \noindent
  \begin{minipage}[t]{0.6\linewidth}
    {\small\bfseries Matem\'aticas Especiales --- Taller (quiz) Parcial I}\\
    {\small Semestre 02 de 2019}
  \end{minipage}%
  \hfill
  \begin{minipage}[t]{0.35\linewidth}
    \raggedleft\small Escuela de Matem\'aticas. Universidad Nacional de Colombia, Sede Medell\'in.
  \end{minipage}\\[2pt]
  \rule{\linewidth}{0.6pt}
}
\newcommand{\examfooter}{%
  \vfill
  \rule{\linewidth}{0.4pt}\\[1pt]
  {\scriptsize\textit{Transcripci\'on digital --- MathUNAL}\hfill\textcolor{unalblue}{\textbf{\thepage}}}
}
\pagestyle{empty}

\begin{document}

\examheader
\vspace{6pt}

\textbf{1.} La gr\'afica del tercer modo fundamental $u_3(x,t)$ que corresponde a la soluci\'on $u(x,t)$ de una cuerda en vibraci\'on de longitud $L=1$ y densidad $\rho$, la cual se encuentra atada en sus extremos, sujeta a una tensi\'on fija $T$, y que comienza a vibrar con velocidad cero, corresponde a:

\begin{center}
\begin{tikzpicture}[scale=0.85]
\foreach \row/\fx/\lbl in {0/{sin(2*pi*\x r)}/a, 1/{-sin(4*pi*\x r)}/b, 2/{cos(3*pi*\x r)}/c, 3/{sin(3*pi*\x r)}/d}{
  \begin{scope}[yshift=-4.6*\row cm]
    \draw[-Latex, thin, gray] (-0.3,0) -- (5.6,0);
    \draw[-Latex, thin, gray] (0,-1.6) -- (0,1.6);
    \foreach \xt in {0.2,0.4,0.6,0.8,1.0} \draw[gray] ({5*\xt},0.05) -- ({5*\xt},-0.05) node[below, black, scale=0.65, yshift=-1pt] {\xt};
    \foreach \yt in {-1.0,-0.5,0.5,1.0} \draw[gray] (0.05,\yt) -- (-0.05,\yt) node[left, black, scale=0.65] {\yt};
    \draw[thick, blue!30!black, domain=0:1, samples=150, variable=\x]
      plot ({5*\x}, {\fx});
    \node[left] at (-1.3,0) {\textbf{(\lbl)}};
  \end{scope}
}
\end{tikzpicture}
\end{center}

\vspace{4pt}
\textit{V\'ease la opci\'on (d) como respuesta indicada en el material fuente (corresponde a \\
$u_3(x,t)=\sin(3\pi x)\cos(3\pi c\,t)$, coherente con la definici\'on del tercer modo fundamental para $L=1$).}

\clearpage
\examheader
\vspace{6pt}

\textbf{2.} Una cuerda de longitud $L=0.5$ m y constante $c=440$ m/s vibra de tal suerte que la frecuencia de su modo fundamental es $f=440$ Hz. Las notas del piano que corresponden aproximadamente las frecuencias de su tercero y quinto modo fundamentales son:
\begin{opciones}
  \item Sol, si
  \item Mi, do\#
  \item Mi, Sol
  \item Sol, do
\end{opciones}
\textit{Respuesta indicada: Mi, do\#.}

\vspace{10pt}
\textbf{3.} Sea $f(x)$ como en el problema anterior ($f(x)=x^2$, $-\pi\leq x\leq \pi$, y $f$ se extiende de manera peri\'odica a toda la recta real). La suma de la serie alternante
\[
s = \frac{1}{1^2} - \frac{1}{2^2} + \frac{1}{3^2} - \frac{1}{4^2} + \cdots
\]
es igual a:
\begin{opciones}
  \item $\dfrac{\pi}{4}$
  \item $\dfrac{\pi^2}{12}$
  \item $\dfrac{\pi^2}{6}$
  \item Por ser alternante, la serie $s$ no converge.
\end{opciones}
\textit{Respuesta indicada: $\pi^2/12$.}

\vspace{10pt}
\textbf{4.} Una cuerda de densidad $\rho = 0.5103\times 10^{-3}$ kg/m y longitud un metro se ata a una pared, en uno de sus extremos. \textquestiondown Cu\'anto debe ser la masa de un objeto que atado al extremo opuesto logre la tensi\'on apropiada para que la cuerda al vibrar produzca una frecuencia de 880 Hertz?
\begin{opciones}
  \item $12.5$ kg
  \item $40.2$ kg
  \item $39.51$ kg
  \item $387.198$ kg
\end{opciones}
\textit{Nota del transcriptor: el material fuente presenta un error de exportaci\'on de Moodle en esta pregunta (las cuatro opciones aparecen marcadas simult\'aneamente como ``correctas''), lo cual es matem\'aticamente imposible para una pregunta de \'unica respuesta. Se omite la indicaci\'on de respuesta correcta a la espera de verificaci\'on manual.}

\clearpage
\examheader
\vspace{6pt}

\textbf{5.} Suponga que $f_1(x), f_2(x), \ldots, f_n(x), \ldots$ son funciones ortogonales definidas en $[a,b]$, bajo el producto ``punto'' $\langle f,g\rangle = \int_a^b f(x)g(x)\,dx$. Suponga que la funci\'on $g(x)$ admite un desarrollo en serie $g(x) = \sum_{k=1}^{\infty} a_k f_k(x)$. Entonces, el coeficiente $a_n$ puede calcularse como:
\begin{opciones}
  \item $a_n = \langle f_n, g\rangle$
  \item $a_n = \langle f_n, f_n\rangle$
  \item $a_n = \dfrac{\langle f_n,g\rangle}{\langle f_n,f_n\rangle}$
  \item Ninguna de las anteriores
\end{opciones}
\textit{Respuesta indicada: $a_n = \dfrac{\langle f_n,g\rangle}{\langle f_n,f_n\rangle}$.}

\vspace{10pt}
\textbf{6.} Las frecuencias $f_n$ de las notas de la escala pitag\'orica que corresponden a las potencias $(3/2)^n$, $n=0,1,2,3$, comenzando en la nota $f_0 = $ La, y de tal manera que no se extiendan m\'as all\'a de una octava son:
\begin{opciones}
  \item $f_0,\ \tfrac{3}{2}f_0,\ \tfrac{9}{4}f_0,\ \tfrac{27}{8}f_0$
  \item $f_0,\ 1.5\,f_0,\ \tfrac{9}{8}f_0,\ \tfrac{27}{16}f_0$
  \item La, la\#, si, do
  \item La, si, do, re
\end{opciones}
\textit{Respuesta indicada: $f_0,\ 1.5\,f_0,\ 9/8\,f_0,\ 27/16\,f_0$.}

\vspace{10pt}
\textbf{7.} Sea
\[
h(x) = \begin{cases} 2, & \text{si } 0\leq x < L \\ -1, & \text{si } -L\leq x < 0 \end{cases}
\]
la cual se extiende de manera peri\'odica a todo $\mathbb{R}$, con periodo $2L$. Sea $s(x)$ la funci\'on determinada por su serie de Fourier. El valor de $s(x)$ en $x=2L$ es:
\begin{opciones}
  \item $1/2$
  \item $2$
  \item No se puede determinar porque $h(x)$ no es continua en ese punto
  \item $-1$
\end{opciones}
\textit{Respuesta indicada: $1/2$.}

\clearpage
\examheader
\vspace{6pt}

\textbf{8.} Suponga que sobre una cuerda de longitud $L$, densidad $\rho$, atada en los extremos, y sometida a una tensi\'on constante $T$, act\'ua en el instante $t$ una fuerza externa vertical determinada por una \textit{funci\'on de densidad de fuerza} $\kappa(x,t)$. Esto significa que sobre un tramo cualquiera de cuerda comprendido entre $x$ y $x+h$ act\'ua una fuerza $F$ en la direcci\'on vertical negativa cuya magnitud est\'a dada por $F = \int_x^{x+h} \kappa(s,t)\,ds$. Su ecuaci\'on de movimiento est\'a dada por:
\begin{opciones}
  \item $\dfrac{\partial^2 u}{\partial x^2} = \dfrac{\rho}{c^2}\left(\dfrac{\partial^2 u}{\partial t^2} + \kappa(x,t)\right)$, \quad $c=\sqrt{T/\rho}$
  \item $\dfrac{\partial^2 u}{\partial x^2} = \dfrac{\rho}{c^2}\left(\dfrac{\partial^2 u}{\partial t^2} - \kappa(x,t)\right)$, \quad $c=\sqrt{T/\rho}$
  \item $\dfrac{\partial^2 u}{\partial x^2} = \dfrac{1}{c^2}\left(\dfrac{\partial^2 u}{\partial t^2} + \dfrac{1}{\rho}\kappa(x,t)\right)$, \quad $c=\sqrt{T/\rho}$
  \item $\dfrac{\partial^2 u}{\partial^2 x} = \dfrac{1}{c^2}\left(\dfrac{\partial u}{\partial t} + \dfrac{1}{g}\right)$, \quad $g=9.8\text{ m/s}^2$
\end{opciones}
\textit{Respuesta indicada: opci\'on (c).}

\vspace{10pt}
\textbf{9.} Las siguiente pregunta hace referencia a una cuerda de longitud $L$ y densidad $\rho$, la cual se encuentra atada en sus extremos, sujeta a una tensi\'on fija $T$. En el instante $t=0$, la cuerda toma la forma de la funci\'on $f(x)$. Denotemos por $u(x,t)$ la funci\'on que representa la forma que asume la cuerda en el instante $t$. El hecho de que la cuerda est\'e atada en sus extremos fuerza a que una soluci\'on no constante de la forma $u(x,t) = F(x)G(t)$ satisfaga que:
\begin{opciones}
  \item $F''(x) - kF(x) = 0$, con $k = -\left(\dfrac{n\pi}{L}\right)^2$
  \item $F''(x) - kF(x) = 0$, con $k = \left(\dfrac{n\pi}{L}\right)^2$
  \item $F''(x) - kF(x) = 0$, con $k>0$
  \item Ninguna de las anteriores.
\end{opciones}
\textit{Respuesta indicada: opci\'on (a).}

\clearpage
\examheader
\vspace{6pt}

\textbf{10.} \textquestiondown Cu\'al de las siguientes funciones, definidas inicialmente en $[-1,1]$, y las cuales se extienden de manera peri\'odica a toda la recta real, sabemos con seguridad admite una serie de Fourier convergente en cada punto de la recta real?
\begin{opciones}
  \item $f(x) = 1/\sqrt{|x|}$, si $x$ es diferente de cero, y $f(0)=0$.
  \item $f(x) = \sqrt{|x|}$
  \item $f(x) = x\sin(1/x)$, si $x\neq 0$ y $f(0)=0$.
  \item Ninguna de las anteriores.
\end{opciones}
\textit{Respuesta indicada: Ninguna de las anteriores.}

\vspace{10pt}
\textbf{11.} Las siguiente pregunta hace referencia a una cuerda de longitud $L$ y densidad $\rho$, la cual se encuentra atada en sus extremos, sujeta a una tensi\'on fija $T$. En el instante $t=0$ la cuerda toma la forma de la funci\'on $f(x)$. Denotemos por $u(x,t)$ la funci\'on que representa la forma que asume la cuerda en el instante $t$. Entonces:
\begin{opciones}
  \item $\dfrac{\partial^2 u}{\partial x^2} = c^2\dfrac{\partial u}{\partial t}$
  \item $\dfrac{\partial^2 u}{\partial x^2} = \dfrac{1}{c^2}\dfrac{\partial^2 u}{\partial t^2}$
  \item $\dfrac{\partial^2 u}{\partial x^2} = \dfrac{1}{c^2}\dfrac{\partial u}{\partial t}$
  \item $\dfrac{\partial^2 u}{\partial x^2} = \dfrac{T}{\rho}\dfrac{\partial u}{\partial t}$
\end{opciones}
\textit{Respuesta indicada: opci\'on (b).}

\vspace{10pt}
\textbf{12.} Sea $f(x) = x^2$, $-\pi\leq x\leq \pi$. Si $f$ se extiende de manera peri\'odica a toda la recta real, su desarrollo en serie de Fourier est\'a dado por la serie:
\begin{opciones}
  \item $\dfrac{\pi^2}{3} + 4\displaystyle\sum_{n=1}^{\infty} \frac{(-1)^n}{n^2}\cos(nx)$
  \item $\dfrac{\pi^2}{3} + 4\displaystyle\sum_{n=1}^{\infty} \frac{1}{n^2}\cos(nx)$
  \item $\dfrac{\pi^2}{3} + 4\displaystyle\sum_{k=0}^{\infty} \frac{(-1)^k}{(2k+1)^2}$ \quad \textit{[sic: as\'i aparece en la fuente, sin el factor $\cos(\cdot)$]}
  \item $\dfrac{\pi^2}{3} + 4\displaystyle\sum_{k=0}^{\infty} \frac{(-1)^k}{(2k)^2}$ \quad \textit{[sic: \'idem]}
\end{opciones}
\textit{Respuesta indicada: opci\'on (a).}

\examfooter
\end{document}
